% paper6_trials.tex -- Paper 6. Typeset alone with this file, or inside the
% volume through main.tex.
\documentclass[../main.tex]{subfiles}
\begin{document}
\begin{paper}{Designing and Analysing Multi-Site Field Trials for Soil Health Research}{Multi-site field trials}
\label{paper:trials}

\begin{paperabstract}
Field trials in soil health research are usually repeated across sites,
because a treatment that works on one soil may fail on another. This paper
describes the design shared by the trials reported in this volume: three
sites in contrasting climatic zones, randomised complete blocks at each site,
and a common set of soil and crop measurements. It sets out how each site is
analysed on its own, how the sites are combined, and how the interaction
between treatment and site is tested and interpreted. The paper serves as the
methodological reference for the other contributions, and all example values
are synthetic.
\end{paperabstract}
\paperkeywords{field experiments; randomised complete block design; analysis of variance; multi-site trials; treatment by site interaction}

\section{Introduction}

Agricultural field experiments must separate the effect of a treatment from
the natural variation between plots. The randomised complete block design
does so by grouping plots into blocks that are as uniform as possible,
assigning every treatment once in each block, and randomising the order of
treatments within blocks, so that differences between blocks can be removed
from the experimental error \parencite{gomez1984}. It remains the most common
design for trials on farms and research stations alike.

A single site, however, represents a single soil and climate. When the same
experiment is repeated at several sites, the question of interest changes
from whether a treatment has an effect to whether its effect is consistent.
A treatment by site interaction indicates that the ranking or the size of
treatment effects differs between sites, and such interactions are the rule
rather than the exception in soil management research, as the variable
response of crop yield to reduced tillage across climates illustrates
\parencite{pittelkow2015}. Experiments conducted over space and time require
their own methods of combined analysis \parencite{gomez1984}.

Soil health trials add a further difficulty, because the response is often
described not by one variable but by several indicators that must be
interpreted together. Converting indicators to a common score before they are
combined, as in the Soil Management Assessment Framework
\parencite{andrews2004}, makes results comparable across sites, and a clear
statement of how indicators were chosen and interpreted is needed if results
are to be used by land managers and policy makers \parencite{bunemann2018}.

This paper describes the design and analysis used across the volume. Its aims
are to document the layout of the trials, to state the statistical model used
at each site and across sites, and to set out the rules by which treatment
effects are reported, so that the results of the other papers can be read on
a common basis.

\input{\subfix{methods/paper6_methods}}

\end{paper}
\end{document}
